\documentclass[11pt,a4paper]{article}
\usepackage[margin=2.2cm,top=2.0cm,bottom=2.2cm]{geometry}
\usepackage{amsmath,amssymb}
\usepackage[T1]{fontenc}
\usepackage{mathptmx}
\usepackage{xcolor}
\usepackage{enumitem}
\usepackage{fancyhdr}
\usepackage{tabularx}
\usepackage{mdframed}
\usepackage{multicol}

\definecolor{navy}{HTML}{1F4E79}
\definecolor{rule}{HTML}{333333}
\definecolor{boxbg}{HTML}{F4F6F9}
\definecolor{keyred}{HTML}{B02020}

% \ANSWERStrue for the tutor key
\newif\ifANSWERS
\input{mode}

\pagestyle{fancy}\fancyhf{}
\renewcommand{\headrulewidth}{0pt}
\renewcommand{\footrulewidth}{0.4pt}
\fancyfoot[L]{\footnotesize\color{navy}SAT Math: Foundations to Mastery \ $\cdot$\ Lines\ifANSWERS\ $\cdot$\ \textbf{\color{keyred}ANSWER KEY}\fi}
\fancyfoot[R]{\footnotesize\color{navy}Page \thepage\ of \pageref{LastPage}}
\usepackage{lastpage}

\setlength{\parindent}{0pt}
\setlength{\parskip}{4pt}

\newcommand{\PartHead}[2]{%
  \vspace{6pt}
  {\color{navy}\rule{\textwidth}{0.8pt}}\\[2pt]
  {\large\bfseries\color{navy}#1}\hfill{\small\itshape #2}\\[-4pt]
  {\color{navy}\rule{\textwidth}{0.4pt}}
  \vspace{4pt}
}
\newcommand{\sitting}[1]{{\footnotesize\textsc{\color{navy}#1}}\par\vspace{2pt}}
\newcommand{\spr}{\\[3pt]\hspace*{6pt}Answer: \underline{\hspace{3.2cm}}\hfill{\footnotesize\itshape(write your answer)}}
\ifANSWERS
  \newcommand{\ans}[1]{\ {\color{keyred}\bfseries #1}}
  \newcommand{\note}[1]{\par{\footnotesize\color{keyred}#1}}
\else
  \newcommand{\ans}[1]{}
  \newcommand{\note}[1]{}
\fi

\begin{document}
\label{firstpage}

\begin{center}
{\LARGE\bfseries\color{navy}Lines}\\[4pt]
{\color{navy}SAT Math $\cdot$ Foundations to Mastery $\cdot$ Week of 4 October}
\ifANSWERS\\[4pt]{\color{keyred}\bfseries ANSWER KEY --- tutor copy}\fi
\end{center}
\vspace{2pt}
{\color{rule}\rule{\textwidth}{1pt}}
\vspace{6pt}

\ifANSWERS\else
Name: \underline{\hspace{6cm}} \hfill Due: \underline{\hspace{3.2cm}}
\vspace{6pt}
\fi

\begin{mdframed}[backgroundcolor=boxbg,linecolor=navy,linewidth=0.6pt,innertopmargin=6pt,innerbottommargin=8pt]
{\bfseries\color{navy}How to use this sheet}\par\vspace{2pt}
This is four short sittings, not one long one. Four half-hours across the week will
do more for you than one long Saturday, and you are far more likely to finish it.

\begin{itemize}[leftmargin=14pt,itemsep=1pt,topsep=3pt]
  \item \textbf{Part A is no calculator.} None. It is 15 minutes and it keeps last week's sign work alive.
  \item \textbf{Part D is the part that pays.} It asks what a number \emph{means}. Read each stem twice before you answer.
  \item When you get one wrong, write it in your error log: what you did, and what you should have done.
\end{itemize}
\end{mdframed}

\vspace{4pt}
\begin{center}
\begin{tabularx}{\textwidth}{@{}l l X r@{}}
\textbf{Sitting} & \textbf{Part} & \textbf{What it drills} & \textbf{Minutes}\\[2pt]
\hline\rule{0pt}{11pt}
Sunday   & A    & Sign retention, no calculator                     & 15\\
Tuesday  & B    & Solving linear equations                          & 25\\
Thursday & C    & Slope, four ways                                  & 25\\
Saturday & D, E & What the number means; six mixed, timed           & 25\\
\hline
\end{tabularx}
\end{center}

% ------------------------------------------------ PART A
\PartHead{Part A --- Sign retention}{No calculator. 15 minutes.}
\sitting{Sunday}
Last week you fixed this. Keep it fixed. \textbf{Write every substitution in brackets}
before you work anything out --- that single habit is what stopped the bleeding.

\begin{enumerate}[leftmargin=20pt,itemsep=7pt]
\item $8-(-5)=$ \ifANSWERS\else\underline{\hspace{3.5cm}}\fi \ans{13}
\item $(-7)^2=$ \ifANSWERS\else\underline{\hspace{3.5cm}}\fi \ans{49}
\item $-7^2=$ \ifANSWERS\else\underline{\hspace{3.5cm}}\fi \ans{$-49$ \ (the minus sits outside the power)}
\item $-6-(-6)=$ \ifANSWERS\else\underline{\hspace{3.5cm}}\fi \ans{0}
\item $\dfrac{-24}{-6}-10=$ \ifANSWERS\else\underline{\hspace{3.5cm}}\fi \ans{$-6$ \ ($4-10$)}
\item Expand: $-4(x-3)$ \ifANSWERS\ans{$=-4x+12$}\else\ \underline{\hspace{5.4cm}}\fi
\item Expand and simplify: $-(5x-9)+2x$ \ifANSWERS\ans{$=-3x+9$}\else\ \underline{\hspace{4cm}}\fi
\item $f(x)=x^2-3x$. \ Find $f(-4)$. \ifANSWERS\ans{$=(-4)^2-3(-4)=16+12=28$}\else\ \underline{\hspace{4cm}}\fi
\item $g(x)=5-2x^2$. \ Find $g(-3)$. \ifANSWERS\ans{$=5-2(9)=-13$}\else\ \underline{\hspace{4cm}}\fi
\item Solve $4-3x>13$. \ifANSWERS\ans{$x<-3$. \ $-3x>9$, then divide by $-3$ \emph{and flip}. Check: $x=-4$ gives $16>13$, true.}\else\ \underline{\hspace{4.4cm}}\fi
\end{enumerate}

\ifANSWERS
\note{\textbf{Marking rule: below 8 of 10 means last week has decayed.} Open the next session with a fresh Part A before touching anything new. Items 2, 3 and 10 are the diagnostic ones.}
\fi

\newpage
% ------------------------------------------------ PART B
\PartHead{Part B --- Solving linear equations}{Calculator allowed. 25 minutes.}
\sitting{Tuesday}
One move per line, and write what you did in the margin. When the variable disappears
from both sides, stop and read what is left: a false statement means no solution, a
true statement means every number works.

\begin{enumerate}[leftmargin=20pt,itemsep=9pt]
\item $5(x-3)=2x+9$. What is the value of $x$? \spr
  \ans{\\ $8$. \ $5x-15=2x+9$, so $3x=24$ and $x=8$.}

\item $7x-4=3x+20$. What is the value of $x$? \spr
  \ans{\\ $6$. \ $4x=24$, so $x=6$.}

\item $2(3x+1)=4(x-5)$. What is the value of $x$? \spr
  \ans{\\ $-11$. \ $6x+2=4x-20$, so $2x=-22$ and $x=-11$.}

\item $-2(x-4)=3x+13$. What is the value of $x$? \spr
  \ans{\\ $-1$. \ $-2x+8=3x+13$, so $-5x=5$ and $x=-1$. Check: $-2(-5)=10$ and $-3+13=10$.}

\item $\dfrac{x}{3}+\dfrac{x}{4}=14$. What is the value of $x$? \spr
  \ans{\\ $24$. \ Multiply every term by 12: $4x+3x=168$, so $7x=168$ and $x=24$.}

\item $\dfrac{2x-1}{5}=\dfrac{x+4}{3}$. What is the value of $x$? \spr
  \ans{\\ $23$. \ Cross-multiply: $3(2x-1)=5(x+4)$, so $6x-3=5x+20$ and $x=23$.}

\item $\dfrac{3}{4}x-2=\dfrac{1}{2}x+1$. What is the value of $x$? \spr
  \ans{\\ $12$. \ Multiply by 4: $3x-8=2x+4$, so $x=12$.}

\item The perimeter of a rectangle is $P=2\ell+2w$, where $\ell$ is the length and $w$ is
  the width. Which equation gives $w$ in terms of $P$ and $\ell$?\\
  \makebox[4.2cm][l]{A) $w=P-2\ell$}\makebox[4.2cm][l]{B) $w=\dfrac{P-2\ell}{2}$}%
  \makebox[4.2cm][l]{C) $w=\dfrac{P}{2}-2\ell$}\makebox[3cm][l]{D) $w=\dfrac{P-\ell}{2}$}
  \ans{\\ B. Subtract $2\ell$, then divide the \emph{whole} side by 2. Check with $P=30$,
  $\ell=4$: $w=11$, and $2(4)+2(11)=30$. A stops before dividing. C divides only the
  first term. D forgets the 2 on $\ell$.}

\item How many solutions does $6x+5=2(3x+1)+3$ have?\\
  \makebox[4.2cm][l]{A) Zero}\makebox[4.2cm][l]{B) Exactly one}%
  \makebox[4.2cm][l]{C) Exactly two}\makebox[3cm][l]{D) Infinitely many}
  \ans{\\ D. The right side is $6x+5$, so the equation reads $6x+5=6x+5$ --- true for
  every $x$. A is the error of treating any vanishing variable as ``no solution''.}

\item How many solutions does $4(2x-3)=8x+7$ have?\\
  \makebox[4.2cm][l]{A) Zero}\makebox[4.2cm][l]{B) Exactly one}%
  \makebox[4.2cm][l]{C) Exactly two}\makebox[3cm][l]{D) Infinitely many}
  \ans{\\ A. $8x-12=8x+7$ leaves $-12=7$, which is false, so nothing works. D is the
  mirror-image error of item 9: same vanishing variable, opposite conclusion.}
\end{enumerate}

\newpage
% ------------------------------------------------ PART C
\PartHead{Part C --- Slope, four ways}{Calculator allowed. 25 minutes.}
\sitting{Thursday}
Slope is one idea wearing four costumes: two points, a table, a graph or equation, and
a line you have to write yourself. Always subtract in the same order top and bottom.

\begin{enumerate}[leftmargin=20pt,itemsep=9pt]
\item A line passes through $(-3,-5)$ and $(1,7)$. What is its slope? \spr
  \ans{\\ $3$. \ $\dfrac{7-(-5)}{1-(-3)}=\dfrac{12}{4}=3$.}

\item A line passes through $(2,-1)$ and $(-4,8)$. What is its slope? \spr
  \ans{\\ $-3/2$. \ $\dfrac{8-(-1)}{-4-2}=\dfrac{9}{-6}=-\dfrac{3}{2}$.}

\item The table shows four points on a line.
  \begin{center}
  \begin{tabularx}{0.62\textwidth}{@{}|X|X|X|X|X|@{}}
  \hline
  $x$ & $-2$ & $1$ & $4$ & $7$\\ \hline
  $y$ & $11$ & $2$ & $-7$ & $-16$\\ \hline
  \end{tabularx}
  \end{center}
  What is the slope of the line?\\
  \makebox[3.4cm][l]{A) $-9$}\makebox[3.4cm][l]{B) $-3$}%
  \makebox[3.4cm][l]{C) $-\frac{1}{3}$}\makebox[3cm][l]{D) $3$}
  \ans{\\ B. $\dfrac{2-11}{1-(-2)}=\dfrac{-9}{3}=-3$. A uses the change in $y$ and never
  divides. C flips the fraction. D drops the negative.}

\item The table shows four points on a line.
  \begin{center}
  \begin{tabularx}{0.62\textwidth}{@{}|X|X|X|X|X|@{}}
  \hline
  $x$ & $0$ & $3$ & $6$ & $9$\\ \hline
  $y$ & $-4$ & $2$ & $8$ & $14$\\ \hline
  \end{tabularx}
  \end{center}
  Which equation describes the line?\\
  \makebox[3.6cm][l]{A) $y=2x-4$}\makebox[3.6cm][l]{B) $y=2x+4$}%
  \makebox[3.6cm][l]{C) $y=6x-4$}\makebox[3.4cm][l]{D) $y=-4x+2$}
  \ans{\\ A. $y$ rises 6 for every 3 in $x$, so the slope is 2, and $x=0$ gives $y=-4$.
  B has the intercept sign wrong. C uses the change in $y$ as the slope. D swaps slope
  and intercept.}

\item In the $xy$-plane, a line passes through $(0,7)$ and falls 2 units for every 5
  units it moves to the right. What is the slope of the line? \spr
  \ans{\\ $-2/5$. \ Falling means negative: $\dfrac{-2}{5}$. The point $(0,7)$ is the
  intercept and does not affect the slope.}

\item A line is given by $4x+3y=9$. What is its slope?\\
  \makebox[3.4cm][l]{A) $4$}\makebox[3.4cm][l]{B) $-\frac{3}{4}$}%
  \makebox[3.4cm][l]{C) $\frac{4}{3}$}\makebox[3cm][l]{D) $-\frac{4}{3}$}
  \ans{\\ D. Rearranged, $y=-\dfrac{4}{3}x+3$. A reads the coefficient of $x$ straight
  off. B flips the fraction. C loses the negative that comes from dividing by $+3$.}

\item A line passes through $(-2,5)$ and has slope $-3$. Which equation describes it?\\
  \makebox[3.8cm][l]{A) $y=-3x+11$}\makebox[3.8cm][l]{B) $y=-3x-1$}%
  \makebox[3.8cm][l]{C) $y=-3x+5$}\makebox[3.4cm][l]{D) $y=-2x+5$}
  \ans{\\ B. $5=-3(-2)+b$ gives $5=6+b$, so $b=-1$. A comes from substituting $+2$
  instead of $-2$. C uses the $y$-coordinate as the intercept. D uses the
  $x$-coordinate as the slope.}

\item A line is given by $3x-2y=12$. What is the $y$-coordinate of its $y$-intercept? \spr
  \ans{\\ $-6$. \ Put $x=0$: $-2y=12$, so $y=-6$. You answered 4 on the diagnostic ---
  that is the $x$-intercept, found by putting $y=0$ instead.}
\end{enumerate}

\ifANSWERS\else\newpage\fi
% ------------------------------------------------ PART D
\PartHead{Part D --- What the number means}{Calculator only where asked. 15 minutes.}
\sitting{Saturday}
Answer these in words, or by choosing a statement. \textbf{Do not calculate} unless an
item tells you to. In every model below, one number is the starting amount and one
number is the rate of change. Decide which is which before you write anything.

This is the highest-value part of the sheet. On the diagnostic you did the arithmetic
correctly and still lost the mark, because the question was asking something else.

\begin{enumerate}[leftmargin=20pt,itemsep=10pt]
\item A repair company charges $C=0.18m+45$ dollars for a visit, where $m$ is the number
  of miles driven to reach the customer. What does the 45 represent?\\
  \underline{\hspace{\linewidth}}
  \ans{\\ The fixed call-out charge in dollars --- what the visit costs before any
  driving, that is, when $m=0$.}

\item For the same model $C=0.18m+45$, what does the 0.18 represent?\\
  \underline{\hspace{\linewidth}}
  \ans{\\ The cost in dollars for each additional mile driven.}

\item A savings balance is modelled by $P=250+12w$ dollars, where $w$ is the number of
  weeks since the account was opened. Which statement is correct?\\
  \makebox[0.47\textwidth][l]{A) The balance grows by \$250 each week.}%
  \makebox[0.47\textwidth][l]{B) The balance grows by \$12 each week.}\\[2pt]
  \makebox[0.47\textwidth][l]{C) The balance reaches \$12 after 250 weeks.}%
  \makebox[0.47\textwidth][l]{D) The account was opened with \$12.}
  \ans{\\ B. The 12 multiplies $w$, so it is the weekly rate; 250 is the opening
  balance. A and D swap the two numbers. C treats both as totals.}

\item The value of a delivery van is modelled by $V=1800-150t$ dollars, where $t$ is the
  number of years since it was bought. Which statement is correct?\\
  \makebox[0.47\textwidth][l]{A) The van loses \$150 of value each year.}%
  \makebox[0.47\textwidth][l]{B) The van gains \$150 of value each year.}\\[2pt]
  \makebox[0.47\textwidth][l]{C) The van was worth \$150 when bought.}%
  \makebox[0.47\textwidth][l]{D) The van loses \$1800 of value each year.}
  \ans{\\ A. The minus sign in front of 150 makes it a loss per year. B ignores the
  sign. C and D swap the starting value and the rate.}

\item A gym membership costs $T=35m+60$ dollars in total, where $m$ is the number of
  months. Work out the total cost of an 8-month membership, then say which part of that
  total does not depend on $m$.\\
  Total: \underline{\hspace{3cm}} \hfill Part that does not depend on $m$:
  \underline{\hspace{3.6cm}}
  \ans{\\ Total $=35(8)+60=280+60=\$340$. The \$60 joining fee does not depend on $m$.
  \emph{This is the diagnostic question.} You found 280 and stopped. The constant is
  part of the total --- always add it back.}

\item A phone plan costs $C=0.05t+19$ dollars a month, where $t$ is the number of text
  messages sent. A student says that doubling the number of texts doubles the monthly
  cost. In one sentence, say why that is wrong.\\
  \underline{\hspace{\linewidth}}
  \ans{\\ Because the \$19 is charged whether or not any texts are sent, so only the
  $0.05t$ part doubles; the total goes up by less than double.}
\end{enumerate}

\ifANSWERS\else\newpage\fi
% ------------------------------------------------ PART E
\PartHead{Part E --- Mixed, timed}{Calculator allowed. 12 minutes, timed.}
\sitting{Saturday}
Set a timer for 12 minutes. Answer every one --- there is no penalty for a wrong answer,
so a blank is a mark thrown away. Then check each answer by substituting it back.

\begin{enumerate}[leftmargin=20pt,itemsep=9pt]
\item What value of $x$ satisfies $4(x+3)=2x+22$?\\
  \makebox[3.4cm][l]{A) $5$}\makebox[3.4cm][l]{B) $9.5$}%
  \makebox[3.4cm][l]{C) $17$}\makebox[3cm][l]{D) $20$}
  \ans{\\ A. $4x+12=2x+22$, so $2x=10$ and $x=5$. B comes from distributing the 4 only
  to $x$. C comes from adding 12 instead of subtracting it. D comes from multiplying
  by 2 at the last step instead of dividing.}

\item A line passes through $(-5,4)$ and $(3,-2)$. What is its slope? \spr
  \ans{\\ $-3/4$. \ $\dfrac{-2-4}{3-(-5)}=\dfrac{-6}{8}=-\dfrac{3}{4}$.}

\item A line is given by $2x-5y=20$. What is the $y$-coordinate of its $y$-intercept?\\
  \makebox[3.4cm][l]{A) $-4$}\makebox[3.4cm][l]{B) $\frac{2}{5}$}%
  \makebox[3.4cm][l]{C) $4$}\makebox[3cm][l]{D) $10$}
  \ans{\\ A. Put $x=0$: $-5y=20$, so $y=-4$. B is the slope. C drops the negative.
  D is the $x$-intercept.}

\item How many solutions does $5(2x-1)=10x+4$ have?\\
  \makebox[4.2cm][l]{A) Zero}\makebox[4.2cm][l]{B) Exactly one}%
  \makebox[4.2cm][l]{C) Exactly two}\makebox[3cm][l]{D) Infinitely many}
  \ans{\\ A. $10x-5=10x+4$ leaves $-5=4$, which is false.}

\item A technician's weekly pay is modelled by $W=4.5h+28$ dollars, where $h$ is the
  number of hours worked. What does the 4.5 represent?\\
  \makebox[0.47\textwidth][l]{A) The pay in dollars for each hour worked.}%
  \makebox[0.47\textwidth][l]{B) The weekly pay before any hours.}\\[2pt]
  \makebox[0.47\textwidth][l]{C) The number of hours worked each week.}%
  \makebox[0.47\textwidth][l]{D) The total weekly pay in dollars.}
  \ans{\\ A. It multiplies $h$, so it is the hourly rate; 28 is the fixed weekly
  amount. B swaps the two numbers. C reads 4.5 as a count of hours. D reads it as a
  total.}

\item A tank is draining. The volume remaining is modelled by $V=500-12t$ litres, where
  $t$ is the number of minutes since draining began. How many litres remain after 15
  minutes? \spr
  \ans{\\ $320$. \ $500-12(15)=500-180=320$. Subtract from the starting 500; do not
  report 180.}
\end{enumerate}

\ifANSWERS\else
\vspace{10pt}
\PartHead{Error log}{Fill this in as you go.}
\begin{center}
\begin{tabularx}{\textwidth}{|c|c|X|X|}
\hline
\textbf{Date} & \textbf{Q} & \textbf{What I did} & \textbf{What I should have done}\\
\hline
\rule{0pt}{22pt} & & &\\ \hline
\rule{0pt}{22pt} & & &\\ \hline
\rule{0pt}{22pt} & & &\\ \hline
\rule{0pt}{22pt} & & &\\ \hline
\rule{0pt}{22pt} & & &\\ \hline
\rule{0pt}{22pt} & & &\\ \hline
\rule{0pt}{22pt} & & &\\ \hline
\end{tabularx}
\end{center}
\vspace{4pt}
\begin{center}\small\itshape
Bring this sheet and your error log to the next session.
\end{center}
\fi

\ifANSWERS
\newpage
\PartHead{Five-minute marking routine}{Tutor}
\begin{enumerate}[leftmargin=20pt,itemsep=4pt]
\item \textbf{Score Part A only.} Below 8/10 means last week has decayed, and the next
  warm-up repairs it before anything new is introduced.
\item \textbf{Read Part D closely.} A wrong answer there is a reading problem, not a
  calculation one --- he can do the arithmetic. Ask him to point to the words in the
  stem that told him which number to name.
\item \textbf{Pick two wrong answers from Part E} and have him narrate his thinking out
  loud. Leave the rest unmarked.
\end{enumerate}
If the error log is filled in, say so out loud. He will do more of what you visibly value.
\fi

\end{document}
